%%%%%%%%%%%%%%%%%%%%%%%%%%%%%%%%%%%%%%%%%%%%%%%%%%%%%%%%%%%%%%%%%
% Contents : The reconciliation chapter
% Translation of the French version
% $Id : 10-grisbi-manuel-planned-en.tex, v 3.0 2026/08 Dominique Brochard
%%%%%%%%%%%%%%%%%%%%%%%%%%%%%%%%%%%%%%%%%%%%%%%%%%%%%%%%%%%%%%%%%


\chapter{Scheduler\label{plannedtransactions}}
%\chapter{Échéancier\label{plannedtransactions}}

The scheduler allows you to plan recurring transactions with specific dates or time intervals. Once a transaction is entered into the scheduler, Grisbi automatically copies the scheduled transaction to the transaction list as soon as its \indexword{due date}\index{due!date} is reached.
%L'échéancier permet de planifier des opérations qui reviennent régulièrement avec des dates ou des intervalles de temps déterminés. Une fois qu'une opération est enregistrée dans l'échéancier, Grisbi recopie automatiquement l'opération planifiée dans la liste des opérations dès que sa \indexword{date d'échéance}\index{date!échéance} est atteinte.


Furthermore, Grisbi displays alerts on the home page when the due dates for these transactions are due (see the section \vref{home-details-homepage}, \menus{Homepage}).
%De plus, Grisbi affiche dans la page d'accueil des alertes au déclenchement des échéances de ces opérations (voir la section \vref{home-details-homepage}, \menus{Affichage de la page d'accueil}).

% espace pour changement de thème

\vspacepdf{2mm}
To access scheduled transactions, click on \menus{Scheduler} in the navigation panel, or select \menus{Scheduled transactions} from the information bar (see the chapter \vref{home}, \menus{Home}).
%Pour accéder aux opérations planifiées, cliquez sur \menus{Échéancier} dans le panneau de navigation, ou sélectionnez \menus{Opérations planifiées} avec la barre d'information (voir le chapitre \vref{home}, \menus{Accueil}).

\begin{figure}[htbp]
	\begin{center}
		\includegraphics[width=0.95\textwidth]{image/screenshot/planned_transactions_display}
	\end{center}
	\caption{Viewing scheduled transactions}
	%\caption{Affichage des opérations planifiées}
	\label{planned-transactions-display-img}
\end{figure}

A calendar appears at the bottom of the navigation panel, and the details panel then contains three elements\refimage{planned-transactions-display-img}:
%Un calendrier s'affiche en bas du panneau de navigation, et le pavé des détails possède alors trois éléments\refimage{planned-transactions-display-img}:
\begin{itemize}
	\item the toolbar;
	%\item la barre d'outils;
	\item the list of scheduled transactions;
	%\item la liste des opérations planifiées;
	\item the transaction/scheduled form.
	%\item le formulaire de saisie des opérations planifiées.
\end{itemize}

% espace pour changement de thème
\vspacepdf{2mm}
You can configure the schedule alert settings in the menu \menus{Edit - Preferences - Generalities - Various settings} (see the section \vref{setup-general-planned}, \menus{Scheduler}).
%Vous pouvez configurer les paramètres des alertes de l'échéancier dans le menu \menus{Éditer - Préférences - Généralités - Paramètres divers} (voir la section \vref{setup-general-planned}, \menus{Échéancier}).


\section{Toolbar\label{plannedtransactions-functions}}
%\section{Barre d'outils\label{plannedtransactions-functions}}

The planned transactions toolbar offers the following functions:
%La barre d'outils des opérations planifiées présente les fonctions suivantes:
\begin{itemize}
	\item \menus{New scheduled}: opens the blank Transaction/Scheduled form;
	%\item \menus{Nouvelle Échéance}: ouvre le formulaire de saisie vierge;
	\item \menus{Delete}: deletes the selected scheduled transaction;
	%\item \menus{Supprimer}: supprime l'opération planifiée sélectionnée;
	\item \menus{Edit}: opens the Transaction/Scheduled form for the selected transaction and allows you to edit it;
	%\item \menus{Éditer}: ouvre le formulaire de saisie pour l'opération sélectionnée et permet de la modifier;
	\item \menus{Notes} or \menus{Frequency/Mode}: displays the \menus{Notes} column instead of the \menus{Frequency} and \menus{Mode} columns, or vice versa;
	%\item \menus{Remarques} ou \menus{Périodicité/Mode}: affiche la colonne \menus{Remarques} à la place des colonnes \menus{Périodicité} et \menus{Mode}, ou inversement;
	\item \menus{Execute}: executes the selected scheduled transaction, i.e. copies it into the list of transactions for the relevant account (section \vref{plannedtransactions-execution});
	%\item \menus{Exécuter}: exécute l'opération planifiée sélectionnée, c'est-à-dire la recopie dans la liste des opérations du compte concerné (section \vref{plannedtransactions-execution});
	\item \menus{View}: a drop-down list allows you to choose how scheduled transactions are displayed: all, according to various frequencies, or according to a custom view (section \vref{plannedtransactions-view}).
	%\item \menus{Affichage}: une liste déroulante permet de choisir l'affichage des opérations planifiées: toutes, selon diverses périodicités, ou selon un réglage personnalisé (section \vref{plannedtransactions-view}).
\end{itemize}


\section{List of scheduled transactions\label{plannedtransactions-list}}
%\section{Liste des opérations planifiées\label{plannedtransactions-list}}


\subsection{Description\label{plannedtransactions-list-description}}
%\subsection{Description\label{plannedtransactions-list-description}}

The list of scheduled transactions is displayed in the details panel\refimage{planned-transactions-list-img}.
%La liste des opérations planifiées s'affiche dans le panneau des détails\refimage{planned-transactions-list-img}.

% image centrée
\begin{figure}[htbp]
	\begin{center}
		\includegraphics[width=0.95\textwidth]{image/screenshot/planned_transactions_list}
	\end{center}
	\caption{List of scheduled transactions}
	%\caption{Liste des opérations planifiées}
	\label{planned-transactions-list-img}
\end{figure}

The column heading bar appears at the top. You can \indexword{widen or narrow a column}\index{column!width} by clicking on the separator between two columns and dragging it. To restore the \indexword{column widths}\index{column!width} to their default values, select the \menus{View - Reset the column width} menu from the menu bar.
%Elle affiche en haut la barre des libellés des colonnes. Vous pouvez \indexword{élargir ou rétrécir une colonne}\index{colonne!largeur} en cliquant sur le séparateur entre deux colonnes et en le déplaçant. Pour rétablir la \indexword{largeur des colonnes}\index{colonne!largeur} à leur valeur par défaut, sélectionnez le menu \menus{Affichage - Réinitialiser la largeur des colonnes} dans la barre de menu.

You can scroll the list of scheduled transactions up or down using the mouse wheel, or by using the mouse and the vertical scroll bar. To scroll left or right, use the mouse and the horizontal scroll bar.
%Vous pouvez déplacer la liste des opérations planifiées vers le haut ou vers le bas avec la molette de la souris, ou bien avec la souris et l'ascenseur vertical. Le déplacement éventuel vers la gauche ou la droite se fait avec la souris et l'ascenseur horizontal.

% espace pour changement de thème
\vspacepdf{2mm}
Each transaction is displayed on a single line and across a maximum of six columns, providing a maximum of six information fields for each scheduled transaction. The display fields, along with the \indexword{column headings}\index{column!heading}, are as follows:
%Chaque opération est affichée sur une seule ligne et au maximum six colonnes, ce qui fait au maximum six champs d'information pour chaque opération planifiée. Les champs d'affichage, ainsi que les \indexword{libellés des colonnes}\index{colonne!libellé}, sont les suivants:
\begin{itemize}
	\item \menus{Date}: the date of the transaction;
	%\item \menus{Date}: la date de l'opération;
	\item \menus{Account}: the account in question;
	%\item \menus{Compte}: le compte concerné;
	\item \menus{Payee}: 
	%\item \menus{Tiers};
	\item \menus{Frequency}: the frequency of the transaction;
	%\item \menus{Périodicité}: la périodicité de l'opération;
	\item \menus{Mode}: the method of execution, automatic or manual;
	%\item \menus{Mode}: le mode d'exécution, automatique ou manuel;
	\item \menus{Amount}: the transaction amount; a \indexword{negative amount}\index{amount!negative} is displayed in \textcolor{red}{red}.
	%\item \menus{Montant}: le montant de l'opération; un \indexword{montant négatif}\index{montant!négatif} est affiché en \textcolor{red}{rouge}.
	% saut de ligne pour indentation correcte de la note dans la liste
	
	\Note{}: if the \menus{Notes} field has been requested, it is displayed in place of the \menus{Frequency} and \menus{Method} fields, resulting in only five columns and five fields being displayed (see section \vref{plannedtransactions-comments}, \menus{Comments on planned transactions}).
	%\Note{}: si l'affichage du champ \menus{Remarques} a été demandé, il s'affiche à la place des champs \menus{Périodicité} et \menus{Mode}, ce qui ne fait plus que cinq colonnes et cinq champs affichés (voir la section \vref{plannedtransactions-comments}, \menus{Remarques des opérations planifiées}).
\end{itemize}

% espace pour changement de thème
\vspacepdf{2mm}
An additional blank line appears just below the last transaction line and is used to create a new scheduled transaction (see section \vref{plannedtransactions-new}, \menus{New Scheduled Transaction}).
%Une ligne supplémentaire vide s'affiche juste en-dessous de la dernière ligne d'opération, et sert à créer une nouvelle opération planifiée (voir la section \vref{plannedtransactions-new}, \menus{Nouvelle opération planifiée}).

% espace pour changement de thème
\vspacepdf{2mm}
To improve the readability of the display, Grisbi alternates between a \colorbox[RGB]{215,215,255}{purple background} and white for each line.
%Pour des raisons de lisibilité de l'affichage, Grisbi présente une alternance de couleurs de \colorbox[RGB]{215,215,255}{fond violet} et blanc à chaque ligne.

% espace pour changement de thème
\vspacepdf{2mm}
By clicking on the \menus{View}\label{plannedtransactions-view} tool in the toolbar, a drop-down list\refimage{planned-transactions-view-select-img} allows you to choose from several display modes for upcoming transactions \textit{from today’s date}:
%En cliquant sur l'outil \menus{Affichage}\label{plannedtransactions-view} dans la barre d'outils, une liste déroulante\refimage{planned-transactions-view-select-img} vous permet de choisir plusieurs modes d'affichage des prochaines occurrences \textit{à compter de la date du jour}:
% espace avant image 5mm

\vspacepdf{5mm}
\noindent
\begin{minipage}{.65\linewidth}
	\begin{itemize}[rightmargin=.6cm]
		\item \menus{Unique view}: the next occurrence of each transaction;
		%\item \menus{Vue unique}: la prochaine occurrence de chaque opération;
		\item \menus{Week view}: due in 7 days;
		%\item \menus{Vue hebdomadaire}: à échéance de 7 jours;
		\item \menus{Month view}: within 30 days;
		%\item \menus{Vue mensuelle}: à échéance de 30 jours;
		\item \menus{Two months view}: within 60 days;
		%\item \menus{Vue bimestrielle}: à échéance de 60 jours;
		\item \menus{Quarter view}: within 90 days;
		%\item \menus{Vue trimestrielle}: à échéance de 90 jours;
		\item \menus{Year view}: within one year;
		%\item \menus{Vue annuelle}: à échéance d'un an;
		\item \menus{Custom view}: within any number of days, weeks, months or years.
		%\item \menus{Vue personnalisée}: à échéance d'un nombre quelconque de jours, semaines, mois ou années.
		% saut de ligne pour indentation correcte de la note dans la liste
		
		\Note{}: the current view is indicated by a tick, and the number of transactions in the selected view is shown in brackets.
		%\Note{}: l'affichage courant y est indiqué par une coche et, entre parenthèses, vous avez le nombre d'opérations dans la vue sélectionnée.
	\end{itemize}
\end{minipage}
\hspace{15pt}
\begin{minipage}{.35\linewidth}
	\vspace{-5pt}
	\centering
	\includegraphics[width=.8\textwidth]{image/screenshot/planned_transactions_view_select}
	%\vspace{-5pt}
	\captionsetup{%						% for a change of options of the caption package
		format=plain,						% to avoid error message with labelsep option below, default=hang
		type=figure,%
		name=Fig.,%							% rename label caption, default="Figure" (or "Table")
		justification=centering,			% centring of text and caption label 
		labelsep=newline}					% define the separator between the label and the caption text
	\caption{View selection}
	%\caption{Choix d'affichage}
	\vspace{-10pt}
	\label{planned-transactions-view-select-img}
\end{minipage}	

\vspacepdf{5mm}
If you select \menus{Custom view}, a new window will appear showing an incremental field to select the number of periods and a drop-down list for the type of period required\refimage{planned-transactions-view-custom-img}.
%Si vous choisissez \menus{Vue personnalisée}, une nouvelle fenêtre affiche un champ incrémental pour choisir le nombre de périodes et une liste déroulante pour le type de la période désirée\refimage{planned-transactions-view-custom-img}.

% image centrée 
\begin{figure}[htbp]
	\begin{center}
		\includegraphics[width=0.5\textwidth]{image/screenshot/planned_transactions_view_custom}
	\end{center}
	\vspace{-10pt}
	\caption{Selecting the custom period view}
	%\caption{Choix de l'affichage de la période personnalisée}
	\vspace{-10pt}
	\label{planned-transactions-view-custom-img}
\end{figure}
% image centrée 


\subsection{Sorts\label{plannedtransactions-list-sorts}}
%\subsection{Tris\label{plannedtransactions-list-sorts}}
%TODO : subsection to update - bug 2391
To change the order in which transactions are displayed, you can \indexword{sort the transactions}\index{sort!transactions} by one of the information fields; proceed as follows:
%Pour changer l'ordre d'affichage des opérations, vous pouvez faire un \indexword{tri sur les opérations}\index{tri!opérations} sur un des champs d'information; procédez comme suit:

\begin{enumerate}
	\item \texttt{right-click} on the column header containing the field selected as the \gls{sorting} criterion;
	%\item cliquez-droit sur le libellé de la colonne qui contient le champ d'information choisi comme critère de \gls{tri};
	\item a context menu displays a list of the fields in that column; select the field you wish to use as the sort criterion;
	%\item un menu contextuel affiche la liste des champs de cette colonne; sélectionnez le champ choisi comme critère de tri;
	\item all transactions are displayed, sorted according to the selected criterion;
	%\item toutes les opérations s'affichent, triées selon le critère choisi;
	\item a small black triangle appears to the right or left of the column heading; it points downwards if the sort is ascending, and upwards if it is descending;
	%\item un petit triangle noir apparaît à droite ou à gauche du libellé de la colonne; il pointe vers le bas si le tri est ascendant, vers le haut si le tri est descendant;
	
	% saut de ligne pour indentation correcte de la note dans la liste
	\Note{}: depending on the theme of your desktop environment or the window manager you are using, these triangles can be replaced with other characters such as +, -, >, <, etc.
	%\Note{}: ces triangles peuvent être remplacés, en fonction du thème de l'environnement de bureau ou du gestionnaire de fenêtres que vous utilisez, par d'autres caractères tels que +, -, >, <, etc.
	\item click on this triangle to change the sort order.
	%\item cliquez sur ce triangle pour changer l'ordre du tri.
\end{enumerate}

\Note{}: sorting the display in this way does not affect the balance for each line, which is calculated after each transaction; depending on the sort order and display options you have chosen, the balance for each line may not appear in chronological order.
%\Note{}: un tri sur l'affichage réalisé ainsi n'affecte en aucun cas le solde à chaque ligne, qui est calculé après chaque opération; suivant le tri et les options d'affichage que vous avez choisi, le solde à chaque ligne peut présenter une évolution non chronologique.

%espace pour changement de thème
\vspacepdf{3mm}
You can also configure the general sort order in the \menus{Edit - Preferences} menu (see the section \vref{setup-operations-list-modes}, \menus{Transaction list behavior}).
%Vous pouvez aussi configurer l'ordre général des tris dans le menu \menus{Éditer - Préférences} (voir la section \vref{setup-operations-list-modes}, \menus{Comportement de la liste des opérations}).

\vspacepdf{3mm}
\Note{}: the choice of the \gls{primary sorting} criterion (transaction date or value date) that may have been selected there will necessarily affect the balance displayed on each transaction line, as Grisbi calculates and displays the balance on each line.
%\Note{}: le choix du critère de \gls{tri primaire} (date d'opération ou date de valeur) qui peut y avoir été fait modifie nécessairement le solde affiché à chaque ligne d'opération, puisque Grisbi calcule et affiche le solde à chaque ligne.




%%%   To change the order in which scheduled transactions are displayed, you can \indexword{sort by transactions}\index{sort!transactions} on one of the information fields; proceed as follows:
%Pour changer l'ordre d'affichage des opérations planifiées, vous pouvez faire un \indexword{tri sur les opérations}\index{tri!opérations} sur un des champs d'information; procédez comme suit:

\begin{enumerate}
	\item click on the column heading containing the information field selected as the \gls{sorting} criterion; the only possible sort criteria are \menus{Date}, \menus{Account} and \menus{Payee};
	%\item cliquez sur le libellé de la colonne qui contient le champ d'information choisi comme critère de \gls{tri}; les seuls critères de tri possibles sont la \menus{Date}, le \menus{Compte} et le \menus{Tiers};
	\item all transactions are displayed, sorted according to the selected criterion;
	%\item toutes les opérations s'affichent, triées selon le critère choisi;
	\item a small black triangle appears to the right of the column heading; it points downwards if the sort is ascending, and upwards if it is descending;
	%\item un petit triangle noir apparaît à droite du libellé de la colonne; il pointe vers le bas si le tri est ascendant, vers le haut si le tri est descendant;
	
	\Note{}: 
	%\Note{}: ces triangles peuvent être remplacés, en fonction du thème de l'environnement de bureau ou du gestionnaire de fenêtres que vous utilisez, par d'autres caractères tels que +, -, >, <, etc.
	
	\item 
	%\item cliquez sur ce triangle pour changer l'ordre du tri.
\end{enumerate}


\section{Form for entering scheduled transactions\label{plannedtransactions-form}}
%\section{Formulaire de saisie des opérations planifiées\label{plannedtransactions-form}}

The scheduled transactions entry form\refimage{planned-transactions-form-img} appears below the list of scheduled transactions. It is similar to the standard transactions entry form, but includes an additional row at the top with four to six fields for specifying:
%Le formulaire de saisie des opérations planifiées\refimage{planned-transactions-form-img} s'affiche en-dessous de la liste des opérations planifiées. Il est similaire au formulaire de saisie des opérations ordinaires, mais il comprend, dans sa partie supérieure, une ligne supplémentaire de quatre à six champs pour indiquer:

% image centrée
\begin{figure}[htbp]
	\begin{center}
		\includegraphics[width=0.95\textwidth]{image/screenshot/planned_transactions_form}
	\end{center}
	\vspace{-10pt}
	\caption{Form for entering scheduled transactions}
	%\caption{Formulaire de saisie des opérations planifiées}
	\label{planned-transactions-form-img}
\end{figure}
% image centrée

\begin{itemize}
	\item the account to be debited; in the case of a transfer between two accounts, you can select either account (source or destination); the transaction type (debit or credit) must be entered in the appropriate field, and the offsetting entry will be automatically created in the other account when the scheduled transaction is executed;
	%\item le compte affecté; dans le cas d'un virement entre deux comptes, vous pouvez sélectionner l'un ou l'autre des deux comptes (origine ou destination); le type de l'opération (débit ou crédit) doit être renseigné dans le champ adéquat, et la contre-opération sera automatiquement créée dans l'autre compte au moment de l'exécution de l'opération planifiée;
	\item the execution method:
	%\item le mode d'exécution:
	\begin{itemize}
		\item \menus{Manual}: you will need to manually execute the entry in the form by selecting the \menus{Execute} function from the function bar or the context menu (see the section \vref{plannedtransactions-execution}, \menus{Executing a scheduled transaction}),
		%\item \menus{Manuel}: vous devrez exécuter manuellement la saisie dans le formulaire en sélectionnant la fonction \menus{Exécuter} dans la barre de fonctions ou dans le menu contextuel (voir la section \vref{plannedtransactions-execution}, \menus{Exécution d'une opération planifiée}),
		\item \menus{Automatic}: the transaction will be automatically entered into the list of transactions on the scheduled date;
		%\item \menus{Automatique}: l'opération sera automatiquement saisie dans la liste des opérations à la date programmée;
	\end{itemize}
	\item the \menus{Frequency}:
	%\item la \menus{Périodicité}:
	\begin{itemize}
		\item \menus{Once},
		%\item \menus{Une fois},
		\item \menus{Weekly},
		%\item \menus{Hebdomadaire},
		\item \menus{Monthly},
		%\item \menus{Mensuelle},
		\item \menus{BimonthlyBimonthly},
		%\item \menus{Bimestrielle},
		\item \menus{Quarterly},
		%\item \menus{Trimestrielle},
		\item \menus{Annually},
		%\item \menus{Annuelle},
		\item \menus{Custom}\refimage{planned-transactions-form-frequency-custom-img}; in this case, you will need to specify in greater detail the number of periods between each entry or request for validation of the transaction, as well as the desired time period (days, weeks, months, years);
		%\item \menus{Personnalisée}\refimage{planned-transactions-form-frequency-custom-img}; dans ce cas, vous devrez préciser plus finement le nombre de périodes entre chaque saisie ou demande de validation de l'opération, ainsi que la période voulue (jours, semaines, mois, années);
		% image centrée
		%\hspace{20pt}
		\item[]	\begin{figure}[htbp]
			\begin{flushright}						% right align
				\captionsetup{%							% for a change of options of the caption package
					justification=centering}			% centring of text and caption label
				\includegraphics[width=0.95\linewidth]{image/screenshot/planned_transactions_form_frequency_custom}
			\end{flushright}
			\caption{Entering a custom frequency}
			%\caption{Saisie d'une périodicité personnalisée}
			\label{planned-transactions-form-frequency-custom-img}
		\end{figure}
		% image centrée
	\end{itemize}
	\item the \menus{Limite date} for due dates; the default selection is empty; double-clicking or pressing the key combination \keys{Ctrl+Enter \return} in this field opens a calendar; you can also enter the date manually using the various key combinations described in the section \vref{transactions-new-dates}, \menus{Entering the date using the keyboard}.
	%\item la \menus{Date limite} des échéances; le choix par défaut est vide; un double-clic ou la combinaison de touches \keys{Ctrl+Entrée \return} dans ce champ ouvre un calendrier; vous pouvez aussi saisir manuellement la date en utilisant les multiples combinaisons de touches telles que décrites dans la section \vref{transactions-new-dates}, \menus{Saisie de la date au clavier}.
\end{itemize}

% espace pour changement de thème
\vspacepdf{2mm}
The lower part of the form is exactly the same as the form for entering ordinary transactions, and is handled in the same way\refimage{planned-transactions-form-img}, see the section \vref{transactions-form}, \menus{Transaction/Scheduled entry form}.
%La partie inférieure du formulaire est strictement identique au formulaire de saisie des opérations ordinaires, et est géré de la même manière\refimage{planned-transactions-form-img}, voir la section \vref{transactions-form}, \menus{Formulaire de saisie des opérations}.

% espace pour changement de thème
\vspacepdf{2mm}
Once the form is displayed, a context menu accessible by \texttt{right-clicking} in an input field allows you to perform the following actions:
%Une fois le formulaire affiché, un menu contextuel accessible par un \texttt{clic-droit} dans un champ de saisie permet d'effectuer les actions suivantes:

\begin{itemize}
	\item \menus{Cut} (the selection);
	%\item \menus{Couper} (la sélection);
	\item \menus{Copy} (the selection);
	%\item \menus{Copier} (la sélection);
	\item \menus{Paste} (the selection);
	%\item \menus{Coller} (la sélection);
	\item \menus{Delete} (the selection);
	%\item \menus{Supprimer} (la sélection);
	\item[]\hspace{-12pt} then
	%\item[]\hspace{-12pt} puis
	\item \menus{Select All} (in the input field);
	%\item \menus{Tout sélectionner} (dans le champ de saisie);
	\item \menus{Insert Emoji}\index{émoji}.
	%\item \menus{Insérer un émoji}\index{émoji}.
\end{itemize}

% espace pour changement de thème
\vspacepdf{2mm}
In addition, another context menu, accessible by right-clicking in the entry form within a grey area outside the input fields, allows you to access the form’s configuration by selecting \menus{Configure the form} (see the section \vref{setup-form}, \menus{Transactions form}).
%De plus, un autre menu contextuel accessible par un clic-droit dans le formulaire de saisie, dans une zone grise en dehors des champs de saisie, permet d'accéder à la configuration du formulaire, en sélectionnant \menus{Configurer le formulaire} (voir la section \vref{setup-form}, \menus{Formulaire des opérations}).

% espace pour changement de thème
\vspacepdf{2mm}
Please note that the choice of transaction types is necessarily linked to the selected account type (see the section \vref{accounts-type}, \menus{Types of Grisbi accounts}).
%Notez que le choix des types d'opérations est nécessairement lié au type du compte sélectionné (voir la section \vref{accounts-type}, \menus{Types de compte de Grisbi}).

%TODO following
\section{Calendrier et prévisions\label{plannedtransactions-calendar}}

Un calendrier\refimage{planned-transactions-calendar-img} s'affiche en bas du panneau de navigation\refimage{planned-transactions-display-img}. Le jour en cours est \colorbox[RGB]{198,200,201}{surligné}, et les dates \textbf{en gras} correspondent aux dates des opérations planifiées de tous les comptes arrivant à échéance.
	
	%Cliquez sur le calendrier pour le sélectionner. Vous pouvez le parcourir avec la souris ou le clavier:
	%\begin{itemize}
	%	\item[]	\begin{Center}
		%			\includegraphics[width=0.65\linewidth]{image/screenshot/planned_calendar}
		%			\end{Center}
	%			\label{planned-transactions-calendar-img}
	%	\item les touches \keys{\arrowkeyleft} ou \keys{\arrowkeyright}, permettent de passer au jour précédent ou suivant;
	%	\item les touches \keys{\arrowkeyup} ou \keys{\arrowkeydown} permettent de passer à la semaine précédente ou suivante;
	%	\item la combinaison de touches \keys{Ctrl+\arrowkeyleft} ou \keys{Ctrl+\arrowkeyright} permet de passer au mois précédent ou suivant;
	%	\item la combinaison de touches \keys{Ctrl+\arrowkeyup} ou \keys{Ctrl+\arrowkeydown} permet de passer à l'année précédente ou suivante.
	%\end{itemize}
	
	\vspace{5mm}
	\noindent
	\begin{minipage}{.6\linewidth}
		\vspace{-15pt}
		\begin{itemize}[rightmargin=.6cm]
			\item les touches \keys{\arrowkeyleft} ou \keys{\arrowkeyright}, permettent de passer au jour précédent ou suivant;
			\item les touches \keys{\arrowkeyup} ou \keys{\arrowkeydown} permettent de passer à la semaine précédente ou suivante;
			\item la combinaison de touches \keys{Ctrl+\arrowkeyleft} ou \keys{Ctrl+\arrowkeyright} permet de passer au mois précédent ou suivant;
			\item la combinaison de touches \keys{Ctrl+\arrowkeyup} ou \keys{Ctrl+\arrowkeydown} permet de passer à l'année précédente ou suivante.
		\end{itemize}
	\end{minipage}
	\hspace{10pt}
	\begin{minipage}{.4\linewidth}
		%\begin{wrapfigure}{R}{0.4\textwidth}%50mm
		%\vspace{-\intextsep}				% space above the floating (minus intextsep=separation between float and text in text)
		\vspace{-10pt}
		%\centering							% centering the floating figure in the "wrap"
		\raggedleft
		\includegraphics[width=.9\textwidth]{image/screenshot/planned_transactions_calendar}
		%\vspace{-15pt}						% space between the floating and the caption below the floating
		\captionsetup{%						% for a change of options of the caption package
			format=plain,					% to avoid error message with labelsep option below, default=hang
			type=figure,%
			name=Fig.,%						% rename label caption, default="Figure" (or "Table")
			justification=centering,		% centring of text and caption label 
			labelsep=newline}				% define the separator between the label and the caption text
		\caption{Calendrier des échéances}		% \caption is mandatory to reference a figure in lof
		%\vspace{-10pt}						% space below the caption
		\label{planned-transactions-calendar-img}
		%\end{wrapfigure}
	\end{minipage}
	
	
	\section{Sélection d'une opération planifiée\label{plannedtransactions-selection}}
	
	Pour sélectionner une opération planifiée, vous avez deux moyens:
	\begin{itemize}
		\item cliquer sur sa ligne;
		\item déplacer la sélection avec les touches \keys{\arrowkeyup} et \keys{\arrowkeydown} du clavier.
	\end{itemize}
	
	L'opération apparaît alors sur \colorbox[RGB]{245,155,155}{fond rose foncé} et ses détails s'affichent dans le formulaire.
	
	\vspace{2mm}
	Un \texttt{clic-droit} sur la ligne d'une opération planifiée affiche un menu contextuel avec les fonctions suivantes\refimage{planned-transactions-contextMenu-img}:
	
	\vspace{3mm}
	\noindent
	\begin{minipage}{.65\linewidth}
		\vspace{-25pt}
		\begin{itemize}[rightmargin=.6cm]
			\item \menus{Éditer l'opération};
			\item \menus{Cloner l'opération};
			\item \menus{Nouvelle opération};
			\item \menus{Supprimer l'opération};
			\item \menus{Afficher les remarques} ou \menus{Afficher Périodicité/Mode};
			\item \menus{Exécuter l'opération}.
		\end{itemize}
	\end{minipage}
	\hspace{10pt}
	\begin{minipage}{.35\linewidth}
		%\begin{wrapfigure}{R}{0.4\textwidth}%50mm
		%\vspace{-\intextsep}				% space above the floating (minus intextsep=separation between float and text in text)
		%\vspace{-20pt}
		\centering							% centering the floating figure in the "wrap"
		%\raggedleft
		\includegraphics[width=.9\textwidth]{image/screenshot/planned_transactions_contextMenu}
		%\vspace{-15pt}						% space between the floating and the caption below the floating
		\captionsetup{%						% for a change of options of the caption package
			format=plain,					% to avoid error message with labelsep option below, default=hang
			type=figure,%
			name=Fig.,%						% rename label caption, default="Figure" (or "Table")
			justification=centering,		% centring of text and caption label 
			labelsep=newline}				% define the separator between the label and the caption text
		\caption{Menu contextuel\\ d'une opération planifiée}		% \caption is mandatory to reference a figure in lof
		%\vspace{-10pt}						% space below the caption
		\label{planned-transactions-contextMenu-img}
		%\end{wrapfigure}
	\end{minipage}
	
	
	\section{Nouvelle opération planifiée\label{plannedtransactions-new}}
	
	La saisie d'une opération planifiée se fait d'une manière similaire à la saisie d'une opération ordinaire.
	
	Pour \indexword{ouvrir ou afficher}\index{formulaire de saisie!afficher} le formulaire de saisie, utilisez une des méthodes suivantes:
	\begin{itemize}
		\item dans la barre de menu, sélectionnez \menus{Éditer - Nouvelle opération} si un compte est sélectionné;
		\item dans la barre de menu, sélectionnez \menus{Affichage} et cochez la case de l'option \menus{Montrer le formulaire de saisie des opérations};
		\item sélectionnez la ligne vide en bas de la liste des opérations et appuyez sur la touche \keys{Entrée \return};
		\item cliquez sur la ligne vide en bas de la liste des opérations;
		\item appuyez sur la combinaison de touches \keys{Ctrl+T};
		\item dans la barre d'outils, cliquez sur \menus{Nouvelle échéance};
		\item dans le menu contextuel de la liste des opérations, sélectionnez \menus{Nouvelle opération};
		\item cliquez sur le petit triangle d'\indexword{affichage/masquage}\index{formulaire de saisie!masquer-afficher} du formulaire de saisie des opérations à gauche de son libellé.
		
		% saut de ligne pour indentation correcte de la note dans la liste
		\vspacepdf{2mm}
		\Note{}: ce triangle peut être remplacé, en fonction du thème de l'environnement de bureau ou du gestionnaire de fenêtres que vous utilisez, par d'autres caractères tels que +, -, >, <, etc.
	\end{itemize}
	
	% espace pour changement de thème
	\vspacepdf{2mm}
	Renseignez les champs avec les paramètres de votre opération planifiée. Pour plus de détails sur la saisie de ces paramètres dans le formulaire, voir les sections \vref{plannedtransactions-form}, \menus{Formulaire de saisie des opérations planifiées} et \vref{plannedtransactions-calendar}, \menus{Calendrier et prévisions}.
	
	
	\section{Ventilation d'une opération planifiée\label{plannedtransactions-breakdown}}
	
	La ventilation d'une opération planifiée se fait d'une manière similaire à la ventilation d'une opération ordinaire. Une fois que vous avez renseigné les champs de la première ligne du formulaire (voir la section \vref{plannedtransactions-form}, \menus{Formulaire de saisie des opérations planifiées}), saisissez les autres paramètres de l'opération dans la partie inférieure de ce formulaire (voir la section \vref{transactions-form}, \menus{Formulaire de saisie des opérations}), et procédez à la saisie de l'opération planifiée ventilée, comme décrit précisément dans la section \vref{transactions-breakdown}, \menus{Ventilation d'une opération}.
	
	
	\section{Modification d'une opération planifiée\label{plannedtransactions-modify}}
	
	Pour modifier une opération planifiée, ou une opération planifiée ventilée ou une de ses sous-opérations, procédez comme suit:
	
	\begin{enumerate}
		\item sélectionnez-la soit:
		\begin{itemize}
			\item par sélection dans la liste des opérations où elle s'affichera sur \colorbox[RGB]{245,155,155}{fond rose foncé},
			\item par \texttt{clic-gauche} sur l'outil \menus{Éditer} de la barre d'outils,
			\item par \texttt{clic-droit} sur la ligne de l'opération puis sélection de \menus{Éditer l'opération} dans le menu contextuel;
		\end{itemize}
		\item le formulaire de saisie s'affiche et le champ \menus{Date} est sélectionné: vous pouvez alors modifier toutes les données de l'opération à votre convenance;
		\item validez par le bouton \menus{Valider}.
	\end{enumerate}
	
	
	\section{Clonage d'une opération planifiée\label{plannedtransactions-duplicate}}
	
	Pour cloner une opération planifiée, sélectionnez-la pour qu'elle s'affiche dans le formulaire de saisie, puis utilisez l'une de ces deux méthodes:
	\begin{itemize}
		\item dans la barre de menus, sélectionnez \menus{Éditer - Cloner l'opération};
		\item dans le menu contextuel, sélectionnez \menus{Cloner l'opération}.
	\end{itemize}
	
	% espace pour changement de thème
	\vspacepdf{2mm}
	L'opération sélectionnée précédemment apparaît toujours sur \colorbox[RGB]{245,155,155}{fond rose foncé}, et l'opération clonée s'affiche dans une nouvelle ligne en-dessous.
	
	Vous pouvez alors éditer cette nouvelle opération pour, par exemple, en modifier la date.
	
	
	\section{Remarques des opérations planifiées\label{plannedtransactions-comments}}
	
	Vous pouvez afficher ou cacher les remarques de deux manières:
	\begin{itemize}
		\item cliquez sur l'outil \menus{Remarques} dans la barre d'outils;
		\item cliquez-droit sur la ligne de l'opération et sélectionnez \menus{Afficher/Masquer les remarques} dans le menu contextuel.
	\end{itemize}
	
	Les remarques de toutes les opérations s'affichent alors à la place des colonnes \menus{Périodicité} et \menus{Mode} dans la liste des opérations planifiées\refimage{planned-transactions-list-comments-img}. Recommencer la même séquence réaffiche les colonnes \menus{Périodicité} et \menus{Mode} à la place de la colonne \menus{Remarques}.
	
	% image centrée
	\begin{figure}[htbp]
		\begin{center}
			\includegraphics[width=0.95\textwidth]{image/screenshot/planned_transactions_list_comments}
		\end{center}
		\caption{Affichage des remarques des opérations planifiées}
		\label{planned-transactions-list-comments-img}
	\end{figure}
	% image centrée
	
	
	\section{Exécution d'une opération planifiée\label{plannedtransactions-execution}}
	
	Cette fonction permet de forcer manuellement l'exécution d'une opération planifiée en dehors de la date prévue.
	
	La fonction \menus{Exécuter} de la barre d'outils sert uniquement à faire la recopie, dans le formulaire de saisie, de l'occurrence d'une opération sélectionnée dans la liste des opérations planifiées. Pour que cette occurrence soit effectivement transférée dans la liste des opérations du compte concerné, il est nécessaire de valider ensuite le formulaire de saisie, après avoir éventuellement modifié l'opération.
	
	% espace pour changement de thème
	\vspacepdf{2mm}
	Pour exécuter une occurrence d'une opération planifiée, procédez comme suit:
	\begin{enumerate}
		\item sélectionnez-la dans la liste;
		\item exécutez-la ensuite avec:
		\begin{itemize}
			\item soit un \texttt{clic-gauche} sur l'outil \menus{Exécuter} de la barre d'outils,
			\item soit un \texttt{clic-droit} sur la ligne de l'opération puis la sélection d' \menus{Exécuter l'opération} dans le menu contextuel;
		\end{itemize}
		\item l'occurrence sélectionnée alors s'affiche dans le formulaire de saisie avec le masquage de la ligne supérieure des paramètres de l'opération planifiée; vous pouvez ensuite la modifier selon votre besoin;
		\item validez avec le bouton \menus{Valider}; l'occurrence de l'opération est transférée dans la liste des opérations du compte concerné, et la liste des opérations planifiées affiche l'occurrence suivante (si elle existe, car l'opération planifiée peut être unique) de l'opération planifiée.
	\end{enumerate}
	
	% espace avant Attention ou Note : 5 mm
	\vspacepdf{2mm}
	\Note{}: après cette validation, le transfert est fait sans autre avertissement.
	
	% espace pour changement de thème
	\vspacepdf{2mm}
	Vous pouvez aussi vérifier que la validation de l'exécution de l'occurrence de l'opération planifiée affiche bien ladite opération dans la liste des opérations du compte concerné.
	
	
	\section{Suppression d'une opération planifiée\label{plannedtransactions-remove}}
	
	Pour supprimer une opération planifiée, procédez comme suit:
	\begin{enumerate}
		\item sélectionnez-la dans la liste;
		\item puis après soit:
		\begin{itemize}
			\item un \texttt{clic-gauche} sur l'outil \menus{Supprimer} dans la barre d'outils,
			\item ou un \texttt{clic-droit} sur la ligne de l'opération planifiée puis la sélection de \menus{Supprimer l'opération} dans le menu contextuel;
		\end{itemize}
		\item une boîte de dialogue s'ouvre\refimage{planned-transactions-delete-img} et vous propose de supprimer:
		\begin{figure}[htbp]
			\begin{center}
				\includegraphics[width=0.5\textwidth]{image/screenshot/planned_transactions_delete}
			\end{center}
			\vspace{-10pt}
			\caption{Suppression d'une opération planifiée}
			\vspace{-10pt}
			\label{planned-transactions-delete-img}
		\end{figure}
		\begin{itemize}
			\item \keys{Toutes les occurrences}: toutes les opérations, courante et futures, seront supprimées, les opérations passées sont conservées, 
			\item \keys{Juste celle-çi}: c'est-à-dire uniquement l'occurrence courante (seule l'opération sélectionnée sera supprimée);
			\item la case à cocher \menus{Gardez ce choix et ne plus voir ce message?} permet de conserver votre choix pour les suppressions futures et la boîte de dialogue ne s'affichera plus;
			
			\vspacepdf{2mm}
			\Attention{}: en cochant cette case, les suppression futures seront immédiates, sans avertissement et sans boîte de dialogue.
			
			\vspacepdf{2mm}
			\Note{}: vous pourrez rétablir l'affichage de la boîte de dialogue de suppression dans les préférence de Grisbi, voir le chapitre \vref{setup-operations-remove} \menus{Messages avant suppression}.
		\end{itemize}
		\item faites votre choix, la suppression est exécutée immédiatement.
	\end{enumerate}
	
	% espace avant Attention ou Note : 5 mm
	\vspacepdf{2mm}
	\Attention{}: immédiatement après le choix entre \menus{Toutes les occurrences} et \menus{Juste celle-ci}, la suppression est exécutée sans autre avertissement.
	
	% espace avant Attention ou Note : 5 mm
	\vspacepdf{2mm}
	\Attention{}: la suppression d'une opération planifiée est \indexword{irréversible}\index{opération!irréversible}!
